% Methods. Every fact is traceable to ../metadata/FACTS_PENDING.md (source codes H, N, V1, V2,
% P, L) or to a script in ../scripts/. \pending{} marks facts still to be supplied.

\section{Methods}

\subsection{Plant material and treatments}
\textit{Arabidopsis thaliana} seedlings of five genotypes were grown on square agar plates in
five lanes per plate: Col-0 (lane A), \textit{fls2-17} (B), \textit{fls2-101} (C),
\textit{cml24-2} (D) and \textit{cml24-4} (E). Lane assignments are taken from the 2016 analysis
folders (FACTS R8). \pending{Allele origins, backgrounds and stock numbers.}
\textit{FLS2} encodes the flagellin receptor that perceives the flg22 epitope
\citep{gomezgomez1999,felix1999,gomezgomez2000,chinchilla2006}. \textit{CML24} encodes a
calmodulin-like Ca$^{2+}$ sensor \citep{delk2005} implicated in flg22-triggered nitric-oxide
signalling \citep{ma2008}.
Seedlings were grown for 7 days on 0.5$\times$ Murashige--Skoog/Linsmaier--Skoog basal salts with
0.1\% (w/v) sucrose and 0.5\,g\,L$^{-1}$ MES (pH 5.7), solidified with 0.8\% Phytagel (FACTS R21)
\pending{light, photoperiod and temperature}.
Plates were treated with a mock solution, flg22 (1\,$\mu$M, the laboratory's standard working
concentration; FACTS R17) or \textit{Pseudomonas syringae} pv.\ tomato DC3000 (FACTS R5, R16)
Bacteria were applied by flood inoculation after \citet{ishiga2011} at an OD$_{600}$ of 0.01 in
10\,mM MgCl$_2$ (FACTS R22). The 2016 protocol draft names Silwet L-77 as the mock
\pending{confirm mock and flg22 application method, FACTS P2}. Four plates per treatment were imaged in the July experiment (set A).

The independent November experiment (set B; FACTS R7, R9) used Col-0, \textit{fls2-17} and
\textit{fls2-101} only, with two or three plates per treatment, three rows per plate and three
seedlings per genotype per row. A pilot experiment (February--March 2016) measured rosette leaf
area from images of eight genotypes, including \textit{glr3.3} and \textit{glr3.4} alleles
(FACTS R10--R11). It is reported in the supplement only.

\subsection{Fluorescence hyperspectral imaging}
Plates were imaged in the Gilroy Laboratory (UW-Madison Botany) with a MacroPhor\texttrademark{}
fluorescence hyperspectral line-scan imager (Middleton Spectral Vision; FACTS R1). It carried an
Andor Zyla~5.5 sCMOS camera in 12-bit low-noise mode (FACTS R2). Set A used 488\,nm excitation
and 500--800\,nm emission (FACTS R3). The camera recorded 324 bands from 400.3 to 808.7\,nm at
about 1.23\,nm spacing. Set B used 405\,nm excitation and 400--800\,nm emission (FACTS R4),
recorded as 406 bands from 400.2 to 800.0\,nm, each row-scan acquired at two excitation laser powers,
50\% and 58\% of maximum (\texttt{Laser50}, \texttt{Laser58}; FACTS R15). Each plate was imaged as three
row-scans, one row of seedlings each.
Set A was imaged at \texttt{Hour0} (11:20--12:26 on 22 July 2016, about 6\,h before the 1\,h scans;
the pre-treatment baseline, FACTS R13) and at 1, 2, 3, 24 and
48\,h after infection (FACTS R6). Set B was imaged at 0, 6, 24 and 48\,h.
Acquisition settings were read from every ENVI header (\texttt{metadata/cube\_manifest.csv}).

\subsection{Data provenance and quality}\label{sec:quality}
All \nFiles{} files on the 2016 drive were inventoried and hashed (\texttt{scripts/00\_inventory.py}).
Every file name was parsed without exception: \nCubes{} hyperspectral cubes (\nCubesA{} set-A
samples, \nCubesB{} set-B samples, plus set-B bacteria-only and reference scans and one reference
and one artefact scan in set A) and \nRoiFiles{} hand-drawn region-of-interest (ROI) spectra
(\nRoiUnique{} unique by content hash). All \nShaRaw{} raw cubes carry SHA-256 checksums in
\texttt{metadata/cube\_manifest.csv}. We found the following problems and corrected for each where possible.
\begin{enumerate}
  \item \textbf{Exposure time changed between timepoints.} It was 0.305\,s at \texttt{Hour0},
  0.203\,s at 1--3\,h, 0.222--0.268\,s at 24\,h (varying within plates) and 0.222\,s at 48\,h.
  The 2016 summary graphs plotted uncorrected intensities, so their apparent drop between 0 and
  1\,h (about 1.5-fold) matches the exposure ratio $0.305/0.203=1.50$. We divide every intensity
  by its own exposure time, and our primary features are exposure-independent ratios
  (Section~\ref{sec:features}).
  \item \textbf{\texttt{Hour0} plates were scanned ``inverted''}, i.e.\ upside down (FACTS R14), except four
  rescans labelled ``NI'' made later the same day at the 1\,h settings. \texttt{Hour0} is used
  only as a baseline and is analysed separately from the upright timepoints.
  \item \textbf{Detector saturation.} At the \texttt{Hour0} exposure, up to
  \satHzeroMaxPct\% of a cube's seedling pixels (median \satHzeroMedPct\%) reached the 12-bit ceiling
  (4095 counts). Saturated pixels
  are excluded from all spectra.
  \item \textbf{Fixed scan order.} Within each timepoint plates were scanned Pathogen, then
  flg22, then Control, 20--40\,min apart, so scan time is partly confounded with treatment. We
  report the drift of Control plates between consecutive hours as a bound on this effect.
  \item \textbf{The plate is the experimental unit.} Treatments were applied per plate (four per
  treatment in set A), so treatment effects are tested at plate level. Genotype contrasts are
  within plates.
  \item \textbf{Inconsistent dark subtraction in the 2016 exports.} \nOffsetRoi{} hand-ROI spectra
  (Pathogen 3\,h plates 1--2 and 24\,h plate 4) were exported with the camera dark level
  (median \offsetMedian{} counts) still in them; all others were dark-subtracted (median
  \offsetCleanMedian{} counts). Because nothing fluoresces at wavelengths shorter than the
  excitation, each spectrum's mean below the excitation line (420--470\,nm in set A, 400--404\,nm
  in set B) was subtracted from it before features were computed (\texttt{scripts/01\_import\_rois.py}).
  \item \textbf{Missing cells.} No hand ROIs were drawn for Pathogen plates at 48\,h, although
  the cubes exist. Our automated segmentation covers them.
\end{enumerate}

\subsection{Spectral processing}\label{sec:features}
Each raw cube was read once, sequentially, and hashed with SHA-256
(\texttt{scripts/02\_harvest\_cubes.py}). The per-band dark level was taken as the median of
background pixels, which sit at about 108 counts. Seedlings were segmented on the dark-subtracted
685\,nm image at 10 background standard deviations, then cleaned, and cotyledon and hypocotyl
components were grouped (\texttt{src/pathofluor/segment.py}). Each segmented seedling was
assigned its lane and plant number from the row layout, which was validated against the 2016
hand ROIs (Supplementary Note~1). Spectral features were computed from the mean unsaturated
spectrum of each seedling after Savitzky--Golay smoothing (\texttt{src/pathofluor/features.py}):
\begin{itemize}
  \item the red/far-red ratio $F_{685}/F_{740}$ \citep{buschmann2007};
  \item the position and full width at half maximum of the red emission peak;
  \item the spectral centroid;
  \item the far-red fraction of 650--800\,nm emission;
  \item the short-wavelength fraction (500--640\,nm emission relative to 650--800\,nm), which
  captures the broad autofluorescence reported by the 2016 vendor analysis;
  \item exposure-normalised intensities (counts\,s$^{-1}$);
  \item per-pixel heterogeneity of the ratio and of the short-wavelength fraction.
\end{itemize}

\subsection{Statistics and classification}
\emph{Hand-ROI features (primary).} For each set, hour and feature we fitted
$y \sim \mathrm{treatment}\times\mathrm{genotype} + (1\,|\,\mathrm{plate})$ with Control and Col-0 as
references (\texttt{statsmodels} MixedLM, REML; features standardised), both on absolute values and
on change from the same plate and genotype at hour~0 (plates Control-4 and flg22-2 excluded from
the latter). $p$-values were adjusted by the Benjamini--Hochberg procedure across all features,
hours and contrasts within a set \citep{benjamini1995} (\texttt{scripts/05\_stats.py}).
\emph{Seedling-level features.} Seedling values were averaged per plate (and per plate and
genotype), expressed as change from hour~0, and compared between treated and Control plates by
Welch's $t$-test, or by the same mixed model for genotype interactions, with the same FDR control
(\texttt{scripts/05c\_seedling\_stats.py}).
\emph{Classification.} Logistic regression (L2, $C=0.1$) and PLS-DA (three components) on the
area-normalised spectrum or on the six shape features were evaluated by leave-one-plate-out
cross-validation. AUC was computed on plate means of out-of-fold scores. Significance came from an
exact permutation test that re-ran the whole cross-validation for every assignment of the treated
label to the same number of plates (\texttt{scripts/06\_models.py}).
\emph{Curve resolution.} MCR-ALS with three non-negative components was fitted to 1,500 unsaturated
seedling pixels per set-A cube (500--800\,nm, dark-subtracted, divided by exposure), initialised
from the purest pixels (\texttt{scripts/05b\_mcr.py}).

\subsection{Data and code availability}
Code (MIT licence), derived data, the 2016 hand-ROI spectra and all manifests are in the
repository. Raw cubes are deposited on Zenodo \pending{DOIs}. The two 2016 Middleton Spectral Vision
analysis reports \citep{msv2016july,msv2017nov} are cited but not redistributed (FACTS R19). The metadata follow MIAPPE where
applicable \citep{papoutsoglou2020} and the FAIR principles \citep{wilkinson2016}.
